\section{Principal opens of general seminorm spectra}\label{sec:general-localization}

Let $A$ be a commutative ring, let $f\in A$, and let $B$ be any localization of
$A$ at the submonoid $\{f^n:n\ge0\}$. No reducedness, finite generation,
or coefficient field is assumed. These statements generalize the principal-open
identification used in \cite[\S5]{FHX}.

\begin{lemma}[Seminorms extend to principal localizations]
\label{spectrum:localization-extension}
\lean{EnvelopeProperty.Berkovich.Spectrum.Away.lift,
EnvelopeProperty.Berkovich.Spectrum.Away.comap_lift,
EnvelopeProperty.Berkovich.Spectrum.Away.comap_injective,
EnvelopeProperty.Berkovich.Spectrum.Away.comap_val_ne_zero,
EnvelopeProperty.Berkovich.Spectrum.Away.val_mk',
EnvelopeProperty.Berkovich.Spectrum.Away.lift_val_mk',
EnvelopeProperty.Berkovich.Spectrum.Away.val_invSelf,
EnvelopeProperty.Berkovich.Spectrum.Away.logAbs_mk'}
\leanok
\uses{spectrum:pullback,spectrum:logarithmic-evaluation}
For $x\in\mathcal S(A)$ with $|f|_x\ne0$, there is a unique seminorm of $B$
restricting to $x$. Its value on a fraction is
$|a/f^n|=|a|_x/|f|_x^n$, and the value on $1/f$ is $|f|_x^{-1}$.
Every seminorm of $B$ restricts to a seminorm with $|f|\ne0$.
The logarithmic fraction formula is
$\log|a/f^n|=\log|a|-n\log|f|$, with $\log|f|$ finite.
\end{lemma}
\begin{proof}
The powers of $f$ lie outside the support of $x$, since their values are the
nonzero numbers $|f|_x^n$. The valuation localization construction therefore
extends $x$ to $B$ and gives the displayed fraction formula. Every element of
$B$ has such a fraction representation, so the formula determines the
extension uniquely. Conversely, the image of $f$ is a unit in $B$ and a
multiplicative seminorm sends a unit to a nonzero value. Evaluating
$f\cdot(1/f)=1$ gives the inverse formula. Applying the extended logarithm
gives the final identity; its denominator has positive finite seminorm.
\end{proof}

\begin{theorem}[Principal-open homeomorphism]\label{spectrum:principal-open}
\lean{EnvelopeProperty.Berkovich.Spectrum.Away.homeomorph,
EnvelopeProperty.Berkovich.Spectrum.Away.range_comap,
EnvelopeProperty.Berkovich.Spectrum.Away.isOpenEmbedding_comap,
EnvelopeProperty.Berkovich.Spectrum.awayHomeomorph,
EnvelopeProperty.Berkovich.Spectrum.isOpenEmbedding_awayComap}
\leanok
\uses{spectrum:localization-extension,spectrum:topology}
Restriction identifies $\mathcal S(B)$ homeomorphically with the open subspace
$\{x\in\mathcal S(A):|f|_x\ne0\}$. In particular it is an open embedding.
The same statement applies to the canonical localization $A_f$.
\end{theorem}
\begin{proof}
Restriction and the extension of the preceding lemma are inverse maps, and
restriction is continuous. For every fraction $a/f^n$, the evaluation of the
extension is the continuous quotient $|a|_x/|f|_x^n$ on this nonvanishing locus.
The coordinate criterion for continuity gives a continuous inverse, hence a
homeomorphism. The locus is open because evaluation at $f$ is continuous and
$\R_{\ge0}\setminus\{0\}$ is open. Identifying with this subspace proves
that restriction is an open embedding. The canonical localization is a
particular instance of the same construction.
\end{proof}

\begin{theorem}[Relative principal opens]\label{spectrum:relative-principal-open}
\lean{EnvelopeProperty.Berkovich.RelativeSpectrum.Away.lift,
EnvelopeProperty.Berkovich.RelativeSpectrum.Away.homeomorph,
EnvelopeProperty.Berkovich.RelativeSpectrum.Away.isOpenEmbedding_comap,
EnvelopeProperty.Berkovich.RelativeSpectrum.awayHomeomorph,
EnvelopeProperty.Berkovich.RelativeSpectrum.isOpenEmbedding_awayComap}
\leanok
\uses{spectrum:principal-open,spectrum:relative}
Let $R$ be a commutative ring with a prescribed real-valued seminorm $v$.
Suppose $A$ and $B$ are $R$-algebras and their structure maps are compatible
with $A\to B$. Then restriction identifies $\mathcal S_v(B)$ homeomorphically
with $\{x\in\mathcal S_v(A):|f|_x\ne0\}$, and is an open embedding.
The same fraction and inverse formulas hold on these relative spectra.
\end{theorem}
\begin{proof}
The absolute extension preserves the value of every element of $A$ and thus
of every scalar in $R$. Compatibility of the structure maps shows that it
preserves the prescribed base seminorm. Restriction has the same property,
so the absolute inverse maps restrict to inverse maps between the relative
spaces. Their evaluation formulas retain continuity on the nonvanishing
relative locus, which is open by continuity of relative evaluation at $f$.
The restricted maps therefore give the asserted homeomorphism and open embedding.
\end{proof}

\begin{corollary}[Trivial-field principal-open specialization {\cite[\S5]{FHX}}]
\label{spectrum:localization-trivial-field}
\lean{EnvelopeProperty.Berk.Away.homeomorph,
EnvelopeProperty.Berk.Away.homeomorph_apply_coe,
EnvelopeProperty.Berk.Away.homeomorph_symm_apply,
EnvelopeProperty.Berk.awayHomeomorph}
\leanok
\uses{spectrum:relative-principal-open,spectrum:trivial-field-bridge}
For an algebra $A$ over a trivially valued field $k$, the principal-open
homeomorphism $(\Spec A_f)^\an\simeq\{x\in(\Spec A)^\an:|f|_x\ne0\}$
is the specialization of the relative-spectrum homeomorphism.
Its forward map is restriction and its inverse is the valuation extension.
\end{corollary}
\begin{proof}
The trivial-field comparison homeomorphism preserves every evaluation and
commutes with restriction. It therefore identifies the nonvanishing subsets
on the two descriptions. Composing this comparison, the relative principal-open
homeomorphism, and the inverse comparison gives the asserted homeomorphism.
Evaluation preservation identifies its forward and inverse maps with the
restriction and valuation extension, respectively.
\end{proof}
